\scanpage{12}
159.

we already have such relations; the idea now is to choose $x,y,z$ so that the relations become as simple as possible. In particular we characterize \uline{all} complete two-dimensional factorial local rings with \uline{algebraically closed} residue field. (Theorem 25.1, Remark 25.2).

THEOREM 25.1. Let $R$ be a two-dimensional local ring with maximal ideal $m$ such that $R/m$ is an algebraically closed field of characteristic $\ne2,3,5$. Assume that $R$ is not regular. The following conditions are equivalent:

(i) The completion $\widehat R$ is factorial.

(i)' $\widehat R$ is normal, and the henselization $R^*$ of $R$ is factorial.

(ii) There exists a basis $\{x,y,z\}$ of $m$ and units $\alpha,\beta$ in $R$ such that
\[
z^2+\alpha y^3+\beta x^5=0.
\]
(iii) There exists a basis $\{x^*,y^*,z^*\}$ of the maximal ideal $m^*$ of $R^*$ such that
\[
(z^*)^2+(y^*)^3+(x^*)^5=0.
\]
(iv) There exists a three-dimensional regular local ring $S$ with regular parameters $u,v,w$ such that
\[
\widehat R\cong S/(u^2+v^3+w^5).
\]
\mspage{11}
